% EGA V, English translation by P. Blass, 1987 (cote 24): archive cote 24, site page 69 (PDF page 70).
% typescript. Draft compared: none.
% AI-transcribed by Claude Opus 5.5 (claude-opus-5-5, adaptive thinking, effort high), Message Batches API, batch ega24b (msgbatch_01RGJJeHdR96U16ZWFuUjtDx), 2026-11-13; not yet checked by eye.
\scanpage{24/69}
\mspage{69}

\uline{Theorem 10.2}

With the above hypothesis the linear pencil $Y_L$ with the
canonical projection $Y_L \rightarrow X$ is X isomorphic in a unique fashion to
the blow-up of prescheme X with center T.

To understand the meaning of this theorem it is convenient to
notice at the beginning of the section that if S = Spec(k) then
for a 'sufficiently general' straight line L in P the condition of
regularity is verified (cf. \uline{catching up indicated in No. 1}) namely
by 5.3.  In what follows in the construction of good linear
pencils.  (\uline{above five lines a mystery}).*
indicated or anticipated at the beginning of the present N°, namely
[illegible]* 5.3.  In what follows in the construction of the
``good'' linear pencils indicated in the beginning of the N°, we
could require that the described pencil should satisfy the said
condition (which is a condition of the same type but different
from the one that consists in requiring that for every $\xi \in L$, $\phi_\xi$
should be $X_{k(\xi)}$ - regular).  We should include the condition in
question in the proposed recapitulating statements.

On the other hand, practically 10.2 is used only in the
situation of 9.5, which makes desirable to announce the
reformulation of 9.5 in terms of pencils which according to 10.2,
in order to be able to include in the statement in question also
the isomorphism of the pencil with a blow up (i.e. to give a
description of the situation permitting a suitable reference).  We
obtain thus a way for every projective smooth geometrically
connected X of dim $\geq 2$ over an infinite field k, to find a closed
smooth subscheme non-empty of codimension two at every one of its

\margin{* Ask AG}
